\chapter{Conclusiones y Trabajos Futuros}
\label{chap:conc}

\section{Conclusiones}

Esta investigación diseñó, implementó y evaluó un prototipo funcional de Sistema de Detección de Intrusiones (IDS) para redes IoT/MQTT, desplegado sobre una arquitectura jerárquica Edge--Fog--Cloud compuesta por microcontroladores ESP32-S3, gateways Raspberry Pi 4 y un servidor PC. El sistema integra modelos compactos TinyML para inferencia en el borde, aprendizaje federado jerárquico (HFL) bidireccional para el entrenamiento colaborativo y cifrado autenticado ligero ASCON-128 para proteger todos los canales de comunicación. La validación experimental se realizó sobre \textbf{ocho combinaciones} (RN/CNN $\times$ ASCON/PLAIN $\times$ NoFOG/FOG) y \textbf{58~intentos} de entrenamiento, consolidados en una capa de observabilidad SRE única que produce métricas, eventos canónicos y trazas comparables a lo largo de todas las variantes (Capítulo~\ref{chap:resu}).

A partir de la evidencia experimental se concluye que el objetivo general fue cumplido: \textbf{el prototipo demuestra que es posible operar un IDS federado, cifrado y embebido sobre hardware IoT real} sin sacrificar de forma significativa la precisión del modelo ni la eficiencia de la comunicación. La mejor variante (\texttt{CNN\_ASCON\_FOG}) alcanzó \textbf{0.975 de accuracy global} con un costo criptográfico inferior al $0.3\%$ del ciclo federado y una huella de memoria de $\sim$4.5\,KB en el ESP32-S3. Más aún, el cifrado autenticado no solo resultó compatible con el ciclo federado, sino que mejoró la convergencia en promedio $+3.2$~puntos, lo que confirma que la triple combinación TinyML--HFL--ASCON puede sostenerse de forma realista en escenarios IoT restringidos.

A continuación se sintetiza una conclusión por cada uno de los cuatro objetivos específicos planteados al inicio del trabajo:

\begin{enumerate}
    \item \textbf{Construcción del pipeline de datos y selección de datasets:} se construyó un pipeline reproducible de captura, normalización y generación de \emph{features} a partir de flujos de red unidireccionales, derivando un conjunto de \textbf{13 características} computables en dispositivos de borde y un problema triclase (\texttt{normal}, \texttt{mqtt\_bruteforce}, \texttt{scan\_A}). La revisión y comparación de datasets públicos relevantes para IDS IoT (CICIoT2023, Edge-IIoTset, TON\_IoT, MQTT-IoT-IDS2020) llevó a fijar MQTT-IoT-IDS2020 como base por su afinidad con el protocolo objetivo, y a documentar las limitaciones de etiquetado heurístico que explican parte del gap residual frente a los modelos centralizados \cite{Neto2023_CICIoT2023, Ferrag2022_EdgeIIoTset, Moustafa2020_TONIoT, Rani2024_IoT_IDS_Datasets}. El pipeline exporta los parámetros de normalización (escalas y medias) en un formato consumido directamente por el firmware del ESP32, eliminando inconsistencias entre el entrenamiento y la inferencia.

    \item \textbf{Comparación y exportación de modelos compactos:} antes de comprometer el ciclo HFL completo se evaluaron en un \emph{notebook} centralizado las arquitecturas candidatas (MLP, CNN-1D, Residual-MLP y un Tiny Transformer), midiendo accuracy, costo computacional y tamaño exportable. El análisis mostró desempeños \emph{offline} cercanos al $90.5\%$ para las tres arquitecturas viables y descartó al Tiny Transformer por exceder los recursos del MCU. Sobre esta base se eligió un MLP/RN de \textbf{1\,139 parámetros ($\approx$4.5\,KB)} como modelo de referencia, con la CNN-1D (con capas Conv y BatchNorm fusionadas y solo las densas finales actualizables) como segunda iteración para evaluar el efecto de mayor capacidad expresiva en el ciclo federado. Tanto los pesos como los parámetros de normalización se exportaron en arreglos C++ embeddable, lo que permitió ejecutar inferencia en el ESP32-S3 con latencias sub-milisegundo y sin asignaciones dinámicas de memoria.

    \item \textbf{Implementación y evaluación del esquema de aprendizaje federado:} se implementó un esquema HFL bidireccional con FedAvg ponderado por número de muestras, soportado por dos topologías intercambiables: una estándar (NoFOG) en la que cada Raspberry Pi entrena y reporta directamente al PC, y una variante con \emph{fog leader/peer} (FOG) en la que un gateway agrega los pesos de sus pares antes de enviar al servidor. El esquema redujo la información intercambiada de $\sim$13.3 millones de valores (todos los flujos) a \textbf{163 parámetros (652\,bytes) por ronda}, lo que representa una reducción del $>99.99\%$ frente al enfoque centralizado y elimina la necesidad de transferir capturas crudas. El comportamiento ante datos non-IID se documentó mediante el desbalance entre gateways (un nodo benigno frente a uno con tráfico mixto) y se cuantificó con métricas de \emph{skew} de accuracy entre gateways. Los hallazgos confirman que la pre-agregación FOG mejora la accuracy entre $+1.5$ y $+2.5$~puntos cuando se combina con cifrado autenticado, pero puede degradar el modelo en su ausencia (\texttt{RN\_PLAIN\_FOG}: $-3.0$~pts), lo que valida empíricamente la necesidad de ASCON como condición habilitante de la capa fog \cite{Imteaj2022_FL_ResourceConstrainedSurvey, Albanbay2025}.

    \item \textbf{Integración del cifrado ligero y validación experimental del sistema completo:} se integró ASCON-128 (estándar NIST 2023~\cite{Kaur2025_ASCON_Survey}) en \emph{todos} los canales del sistema (ESP32~$\leftrightarrow$~RPi, RPi~$\leftrightarrow$~PC y, en el modo FOG, \texttt{peer}~$\leftrightarrow$~\texttt{leader}), garantizando confidencialidad, integridad y autenticidad sin puntos ciegos. La validación experimental sobre las 8 variantes y los 58 intentos cuantificó simultáneamente los tres ejes solicitados: (i)~la \textbf{capacidad de detección} en los nodos, con accuracies entre 0.93 y 0.975 y un techo en \texttt{CNN\_ASCON\_FOG}; (ii)~el \textbf{impacto del esquema federado} frente a centralizar, con un gap acotado a $2$--$7$~puntos según la variante (consistente con la literatura de FL non-IID); y (iii)~los \textbf{efectos en latencia y costo de comunicación}, con un overhead temporal de ASCON entre el $0.10\%$ y el $0.27\%$ del ciclo federado y un overhead de tamaño del $\sim$48\% por mensaje ($\sim$1.07\,KB absolutos), perfectamente compatible con redes IoT típicas. El hallazgo más relevante es que la habilitación de ASCON \emph{no penaliza} la convergencia: las cuatro variantes \texttt{*\_ASCON\_*} promedian $0.9592$ frente a $0.9272$ de las \texttt{*\_PLAIN\_*} ($+3.2$~pts), efecto interpretado como filtrado implícito de mensajes corruptos o duplicados gracias a la verificación obligatoria del tag autenticado.
\end{enumerate}

En síntesis, el conjunto de los cuatro objetivos específicos se cumplió de forma cuantificable y se verificó sobre hardware real, no solo en simulación. La arquitectura propuesta es \textbf{modular} (los tres ejes modelo/seguridad/topología son intercambiables), \textbf{auditable} (cada ronda emite un registro canónico que la capa SRE consolida en \texttt{Analisis de Modelos/Completo/}) y \textbf{viable a escala edge} (el ESP32 conserva $>$60\,\% de SRAM libre durante la operación). El trabajo aporta así evidencia experimental original sobre la triple combinación TinyML--HFL--ASCON, complementando los \emph{surveys} recientes de Kaur~\textit{et~al.}~\cite{Kaur2025_ASCON_Survey} y Dini~\textit{et~al.}~\cite{Dini2024_EmbeddedIIoT_Review} con mediciones \emph{end-to-end} en software embebido, y constituye una base reutilizable para evaluar futuras extensiones sobre el mismo testbed sin necesidad de re-instrumentar la capa de observabilidad.

% ====================================================================
\section{Trabajos Futuros}
\label{sec:future}

Como se delimitó en la Sección~\ref{sec:threat_model} y en las limitaciones del Capítulo~\ref{chap:resu}, el sistema deja por fuera tres familias de amenazas que motivan el siguiente \emph{roadmap}: (i)~clientes federados maliciosos y \emph{data poisoning}, (ii)~ataques de canal lateral sobre la implementación criptográfica y (iii)~inferencia de membresía sobre los pesos federados. A partir de estas limitaciones y de las oportunidades emergentes detectadas durante la evaluación experimental, se proponen las siguientes líneas de investigación:

\begin{itemize}
    \item \textbf{Robustez bizantina y agregación tolerante a clientes maliciosos:} incorporar agregaciones robustas frente a actualizaciones envenenadas (Krum, FLTrust, Fed-NGA) o esquemas jerárquicos como EdgeTrust-Shard~\cite{Tran2026_EdgeTrustShard}, que reportan tolerancia hasta $30\%$ de adversarios manteniendo $\sim 94\%$ de accuracy con apenas 140\,KB de memoria en RPi Pico, complementando los algoritmos avanzados de FL non-IID (FedProx~\cite{Li2020FedProx}, SCAFFOLD, FL personalizado) que reducirían el \emph{gap} de accuracy frente al entrenamiento centralizado siguiendo los hallazgos de \cite{Roy2023_FL_AdaptedModel}.

    \item \textbf{Captura de tráfico real y evaluación de larga duración:} sustituir la simulación de tráfico en los ESP32 por captura y análisis de tráfico MQTT real en un entorno IoT operativo (cámaras, sensores domóticos, instrumentación industrial), eliminando la dependencia de la función heurística de etiquetado en el gateway y midiendo simultáneamente el consumo energético de los nodos durante la inferencia continua y la comunicación cifrada~\cite{Tekin2023_OnDeviceEnergy}, con corridas de semanas que permitan observar la deriva del modelo (\emph{concept drift}) y disparar reentrenamientos federados de forma adaptativa con detectores como ADWIN o Page-Hinkley.

    \item \textbf{Escalabilidad multi-gateway y multi-zona:} extender el despliegue a múltiples Raspberry Pi actuando como gateways independientes con datasets locales distintos, evaluando el impacto de una topología HFL más realista con mayor heterogeneidad y validando empíricamente las ventajas de la pre-agregación FOG en redes con decenas de nodos~\cite{Wang2019_DataFusion_CPSS}, lo que permitiría medir la \emph{escalabilidad horizontal} del esquema más allá del testbed mínimo presentado en este trabajo.

    \item \textbf{Cuantización del modelo y despliegue en MCUs aún más restringidos:} aplicar cuantización post-entrenamiento (INT8) al MLP para reducir aún más la huella de memoria y evaluar su impacto en accuracy, habilitando el despliegue en microcontroladores de menor categoría como ESP32-C3 (400\,KB de SRAM) o en plataformas embebidas industriales~\cite{Dini2024_EmbeddedIIoT_Review}, donde el factor de forma y el consumo son aún más críticos.

    \item \textbf{Endurecimiento criptográfico y comparación con ASCON-128a:} evaluar la variante ASCON-128a (\emph{rate} de 128 bits frente a 64 bits en ASCON-128) para determinar si la mayor tasa de procesamiento reduce el overhead temporal en mensajes de pesos de mayor tamaño, y aplicar las contramedidas frente a ataques de canal lateral (\emph{masking}, \emph{hiding}, \emph{randomized SBox}) reportadas por Kaur~\textit{et~al.}~\cite{Kaur2025_ASCON_Survey}; en paralelo, replicar el patrón experimental de Mirzaali~\textit{et~al.}~\cite{Mirzaali2025_LWC_Permutations} ejecutando ataques activos sobre la Raspberry Pi para validar empíricamente la robustez del esquema frente a adversarios reales.

    \item \textbf{Privacidad diferencial y \emph{transfer learning} on-board:} añadir ruido calibrado a las actualizaciones de pesos (DP-SGD) antes del envío al PC para acotar la inferencia de membresía y robustecer la garantía de privacidad más allá del simple no compartir datos crudos; complementariamente, integrar \emph{transfer learning} local en cada gateway siguiendo la receta de Ficco~\textit{et~al.}~\cite{Ficco2024_FL_TinyML_Onboard} para acelerar la adaptación a tráfico no visto durante el preentrenamiento offline.
\end{itemize}

En conjunto, estas líneas apuntan a transformar el prototipo presentado en una herramienta \emph{operativa} de seguridad IoT lista para escenarios industriales, conservando la capa SRE como soporte de medición y permitiendo comparar de manera reproducible las extensiones futuras frente al \emph{baseline} establecido en este trabajo.
